\icode{N24H-MG-CA-L2-261019-001}
\nceeID[7]{2024上海}{四（2）}
如图所示，太阳系中木星沿着椭圆轨道绕太阳运行，其中 \(a\)、\(c\) 为椭圆长轴上的两个端点，\(b\)、\(d\) 为椭圆短轴上的两个端点。木星绕太阳的运行周期为 \(T\)。

\begin{tikzpicture}[font=\scriptsize]
	\draw (0,0) ellipse (2 and 1.2);
	\draw[fill=red!70] (-1.5,0) circle(.15);
	\foreach \x/\a/\b in {(-2,0)/a/left,(0,-1.2)/b/below,(2,0)/c/right,(0,1.2)/d/above}
		\fill \x node[\b]{\(\a\)} circle(.05);
	\node at(-1.2,-.3)[below]{太阳};
\end{tikzpicture}

（1）某时刻木星沿逆时针绕过 \(a\) 点，从该时刻起，经过 \(\frac{T}{4}\) 后，木星位于\paren[]
\begin{choices}
	\item \(ab\)间
	\item \(b\)点
	\item \(bc\)间
	\item \(c\)点
\end{choices}
（2）已知太阳质量为 \(M\)，引力常量为 \(G\)。若木星在距离太阳 \(r_1\) 处的运行速度大小为 \(v_1\)。求木星运行到距太阳 \(r_2\) 处的速度大小 \(v_2\)。（已知无穷远处引力势能为 0，物体距星球球心的距离为 \(r\) 时的引力势能 \(E_p = -\frac{Gm_{\text{星}}m_{\text{物}}}{r}\)）
\begin{answer}
	（1）C\:（2）\(v_2 = \sqrt{v_1^2 + 2GM(\frac{1}{r_2} - \frac{1}{r_1})}\)\qquad 提示：机械能守恒
\end{answer}